% from ega3-6-fr.tex
\oldpage[III]{139}

où $\mathscr H^q(f,\mathscr K^\bullet)$ est le complexe dont le composant de
degré $i$ est
$\mathscr H^q(f,\mathscr K^i)=\mathrm R^q f_*(\mathscr K^i)$ (loc. cit.).
Rappelons aussi que lorsque $Y$ est réduit à un point, on note
l'hypercohomologie correspondante $\mathbf H^\bullet(X,\mathscr K^\bullet)$
(qui est formée de modules sur $\Gamma(X,\mathscr O_X)$ indépendants du
préschéma ponctuel $Y$ considéré); lorsque $Y=X$ et $f=1_X$, on a
$\mathscr H^n(f,\mathscr K^\bullet)=\mathscr H^n(\mathscr K^\bullet)$
(cohomologie du complexe $\mathscr K^\bullet$); lorsque
$\mathscr K^i=0$, sauf pour $i=i_0$, on a
\[
  \mathscr H^n(f,\mathscr K^\bullet)
    =\mathrm R^{n-i_0}f_*(\mathscr K^{i_0}).
\]
La suite spectrale $'\mathscr E(f,\mathscr K^\bullet)$ est toujours
régulière; les deux suites spectrales sont birégulières lorsque
$\mathscr K^\bullet$ est limité inférieurement (0, 12.4.1).

Tout homotopisme $h:\mathscr K^\bullet\to\mathscr K'^\bullet$ de complexes de
$\mathscr O_X$-Modules (6.1.4) donne un isomorphisme
$\mathscr H^\bullet(f,\mathscr K^\bullet)\xrightarrow{\sim}
\mathscr H^\bullet(f,\mathscr K'^\bullet)$ pour l'hypercohomologie. Il en est
de même lorsqu'on suppose seulement que
$\mathscr H^\bullet(h):\mathscr H^\bullet(\mathscr K^\bullet)\to
\mathscr H^\bullet(\mathscr K'^\bullet)$ est un isomorphisme et que
$\mathscr K^\bullet$ et $\mathscr K'^\bullet$ sont limités inférieurement,
comme il résulte aussitôt de (0, 11.1.5) appliqué à la suite spectrale
(6.2.1.2) et à l'analogue pour $\mathscr K'^\bullet$. Enfin, pour tout
recouvrement ouvert $\mathfrak U=(U_\alpha)$ de $X$, on a aussi défini
(0, 12.4.5) l'hypercohomologie
$\mathbf H^\bullet(\mathfrak U,\mathscr K^\bullet)$ comme la cohomologie du
bicomplexe $C^\bullet(\mathfrak U,\mathscr K^\bullet)$ (dont le composant
d'indices $(i,j)$ est par définition
$C^i(\mathfrak U,\mathscr K^j)$); les
$\mathbf H^n(\mathfrak U,\mathscr K^\bullet)$ sont encore des modules sur
$\Gamma(X,\mathscr O_X)$.
\end{env}

\begin{proposition}[6.2.2]
Soient $X$ un schéma, $\mathfrak U=(U_\alpha)$ un recouvrement de $X$ par des
ouverts affines. Pour tout complexe $\mathscr K^\bullet$ de
$\mathscr O_X$-Modules quasi-cohérents, les modules d'hypercohomologie
$\mathbf H^\bullet(X,\mathscr K^\bullet)$ et
$\mathbf H^\bullet(\mathfrak U,\mathscr K^\bullet)$ sont canoniquement
isomorphes.
\end{proposition}

\begin{proof}
En effet, toute intersection finie $V$ d'ouverts du recouvrement
$\mathfrak U$ est affine (I, 5.5.6), donc
$\mathrm H^q(V,\mathscr K^i)=0$ pour tout $i$ et tout $q>0$ (1.3.1); la
proposition est donc un cas particulier de (0, 12.4.7).
\end{proof}

\begin{proposition}[6.2.3]
Soit $f:X\to Y$ un morphisme quasi-compact et séparé de préschémas. Pour tout
complexe $\mathscr K^\bullet$ de $\mathscr O_X$-Modules quasi-cohérents, les
$\mathscr O_Y$-Modules $\mathscr H^n(f,\mathscr K^\bullet)$ sont
quasi-cohérents.
\end{proposition}

\begin{proof}
Comme les $\mathscr H^q(f,\mathscr K^i)=\mathrm R^q f_*(\mathscr K_i)$ sont
des $\mathscr O_Y$-Modules quasi-cohérents (1.4.10), il en est de même de
$'\mathscr E_2^{pq}$, qui, d'après (6.2.1.1), est quotient d'un noyau
d'homomorphisme de Modules quasi-cohérents par une image d'un tel
homomorphisme (I, 4.1.1). Pour la même raison, tous les
$\mathscr O_Y$-Modules $'\mathscr E_r^{pq}$,
$B_k('\mathscr E_2^{pq})$, $Z_k('\mathscr E_2^{pq})$ de la première suite
spectrale sont quasi-cohérents. La régularité de la suite spectrale
$'\mathscr E(f,\mathscr K^\bullet)$ entraîne que
$Z_\infty('\mathscr E_2^{pq})$ est égal à un des
$Z_k('\mathscr E_2^{pq})$, donc est quasi-cohérent, et il en est de même de
$B_\infty('\mathscr E_2^{pq})=\varinjlim_k B_k('\mathscr E_2^{pq})$
(0, 11.2.4 et I, 4.1.1); les $'\mathscr E_\infty^{pq}$ sont donc aussi
quasi-cohérents. La suite spectrale précédente étant régulière, la filtration
des $F^p(\mathscr H^n(f,\mathscr K^\bullet))$ est discrète et exhaustive;
autrement dit, le $\mathscr O_Y$-Module
$\mathscr H^n(f,\mathscr K^\bullet)$ est réunion d'une suite croissante
$(\mathscr G_k)_{k\geq0}$ de $\mathscr O_Y$-Modules telle que
$\mathscr G_0=0$ et que chaque $\mathscr G_k/\mathscr G_{k-1}$ soit égal à un
des $\mathscr O_Y$-Modules $'\mathscr E_\infty^{pq}$, donc soit
quasi-cohérent. Par récurrence sur $k$, on en déduit que les $\mathscr G_k$
sont quasi-cohérents (1.4.17), et comme
$\mathscr H^n(f,\mathscr K^\bullet)=\varinjlim_k\mathscr G_k$, la proposition
est démontrée (I, 4.1.1).
\end{proof}

